% Feasibility report — replacing the external LLM reasoner.
%
% Build (tectonic fetches packages on demand; no local TeX install needed):
%   tectonic -X compile docs/30-narrator-feasibility-report.tex --outdir docs
%
% The built PDF is committed alongside this source in docs/, matching the other
% tracked PDFs there, so readers who do not run LaTeX still get the report from a
% plain checkout. Tectonic writes only the .pdf (intermediates are skipped), so
% building into docs/ leaves no aux files behind.
%
% Figures are transcribed from:
%   outputs/reasoner_comparison_real/{narrator_arch_sweep,faithfulness_audit}.json
%   outputs/generalizability/generalizability_report.json
%   docs/28-reasoner-cnn-vs-distilledqwen.md, docs/29-narrator-verification-and-gaps.md

\documentclass[11pt,a4paper]{article}

\usepackage[a4paper,margin=24mm,top=25mm,bottom=23mm]{geometry}
\usepackage[T1]{fontenc}
\usepackage{XCharter}
\usepackage{inconsolata}
\usepackage[table]{xcolor}
\usepackage{booktabs,tabularx,array,longtable}
\usepackage{float}
\usepackage{titlesec,enumitem,microtype}
\usepackage{tcolorbox}
\tcbuselibrary{breakable}
\usepackage{fancyhdr}
\usepackage{caption}
\usepackage[colorlinks=true,allcolors=accent]{hyperref}

% ── palette: cool clinical neutrals, teal accent, semantic verdict colours ──
\definecolor{accent}{HTML}{0E7C86}
\definecolor{goGreen}{HTML}{1B7F4B}
\definecolor{holdAmber}{HTML}{8A5A0F}
\definecolor{stopRed}{HTML}{AC3128}
\definecolor{inkMuted}{HTML}{4A5860}
\definecolor{ruleGrey}{HTML}{D3DCE1}
\definecolor{softGrey}{HTML}{F2F6F7}
\definecolor{softTeal}{HTML}{E4F1F2}

\newcolumntype{Y}{>{\raggedright\arraybackslash}X}
\newcolumntype{L}[1]{>{\raggedright\arraybackslash}p{#1}}
\newcommand{\code}[1]{\texttt{\small #1}}
\newcommand{\verd}[2]{\textbf{\textcolor{#1}{#2}}}
\newcommand{\eyebrow}[1]{%
  {\footnotesize\sffamily\bfseries\textcolor{accent}{\MakeUppercase{#1}}}}

\captionsetup[table]{
  font={small,color=inkMuted}, labelfont={bf,color=inkMuted},
  justification=raggedright, singlelinecheck=false, skip=5pt}

\titleformat{\section}{\Large\bfseries}{\thesection.}{0.6em}{}
\titleformat{\subsection}{\normalsize\bfseries}{\thesubsection}{0.6em}{}
\titlespacing*{\section}{0pt}{20pt}{7pt}
\titlespacing*{\subsection}{0pt}{13pt}{4pt}

\pagestyle{fancy}
\fancyhf{}
\renewcommand{\headrulewidth}{0.4pt}
\fancyhead[L]{\footnotesize\sffamily\textcolor{inkMuted}{Feasibility: replacing the external LLM reasoner}}
\fancyhead[R]{\footnotesize\sffamily\textcolor{inkMuted}{\thepage}}

\setlength{\parskip}{6pt}
\setlength{\parindent}{0pt}
\renewcommand{\arraystretch}{1.25}
% a little slack so hyphen-dense lines stretch instead of spilling past the margin
\setlength{\emergencystretch}{2em}

\newtcolorbox{callout}[1][accent]{
  colback=softGrey, colframe=#1, boxrule=0pt, leftrule=2.5pt,
  arc=0pt, outer arc=0pt, left=10pt, right=10pt, top=8pt, bottom=8pt,
  breakable}


% Academic progress report, addressed from the project group to the supervisor.
\begin{document}

\begin{flushleft}
\eyebrow{Project progress report \quad\textbullet\quad RetinalAI screening pipeline}

\vspace{6pt}
{\huge\bfseries Replacing the external LLM reasoner\par}

\vspace{6pt}
{\Large\color{inkMuted} A feasibility and deployment study\par}

\vspace{10pt}
{\footnotesize\sffamily\color{inkMuted}
Revision~4 \textbullet\ 12 August 2026 \quad\textbar\quad
Prepared for the project supervisor \quad\textbar\quad
Teacher model: self-hosted Qwen3-8B-AWQ (vLLM) \quad\textbar\quad
Evaluation data: 240 RFMiD images\par}
\end{flushleft}

\vspace{2pt}
{\color{ruleGrey}\hrule height 1.5pt}
\vspace{10pt}

\section*{Abstract}
\addcontentsline{toc}{section}{Abstract}

We were asked to obtain a lightweight model for the screening pipeline's
reasoning stage, requiring modest inference compute, exceeding 75\% accuracy
with precision at least equal to the incumbent, and to establish
generalisability before any container image was produced. We report that the
requirement is satisfied for the \emph{triage} sub-task by a 3\,KB logistic
regression that attains accuracy and macro-precision of 1.000 at 0.14\,ms per
inference, and that this model has been deployed and verified in service. We
further report that the \emph{narrative} sub-task is not yet suitable for
clinical enablement, and we set out the measurements supporting that judgement.
We draw the supervisor's attention to two findings that qualify the headline
result: the triage task proved to be deterministic, so the perfect scores
indicate an eliminated dependency rather than a modelling advance; and a
quantisation projection made in our previous revision has been tested and found
incorrect, which we correct here.

\section{Objectives}

The brief specified four conditions. We restate them and record where each is
addressed.

\begin{table}[H]
\begin{tabularx}{\textwidth}{@{}L{52mm}L{18mm}Y@{}}
\toprule
\textbf{Condition} & \textbf{Status} & \textbf{Evidence} \\
\midrule
A lightweight final model not requiring substantial inference compute &
\verd{goGreen}{Met} &
3\,KB of JSON weights; 0.14\,ms per inference on CPU; the serving package
depends only on the Python standard library (\S3.1) \\
\addlinespace[3pt]
Alternative architectures to be explored &
\verd{goGreen}{Met} &
Seven candidates evaluated on a common split; five attained 1.000 (\S3.2) \\
\addlinespace[3pt]
Accuracy above 75\% with precision at least equal to the incumbent &
\verd{goGreen}{Met} &
Accuracy 1.000 and macro-precision 1.000, against a threshold of 0.75 and an
incumbent precision of 0.492 (\S3.2) \\
\addlinespace[3pt]
Generalisability established before a container image is produced &
\verd{goGreen}{Met} &
Four independent tests conducted prior to the build (\S3.3) \\
\bottomrule
\end{tabularx}
\end{table}

\begin{callout}[stopRed]
\eyebrow{An important qualification}

\vspace{3pt}
We wish to be explicit that the perfect scores reported above should not be
read as evidence of a superior model. On the out-of-sample set, the triage
decision of the 8B teacher was identical to the disease classifier's own
\code{referral\_priority} field in \textbf{160 of 160 cases}. The task is
therefore deterministic, and the deterministic rules, a 3\,KB linear model and
an 8B language model are indistinguishable on it. The contribution of this work
is the removal of a language-model dependency --- with its associated latency,
operating cost and clinical-data egress --- from the decision path, rather than
an improvement in predictive accuracy.
\end{callout}

\section{Method}

The pipeline delegates two distinct tasks to an external language model, and we
treated them separately throughout, as they are not comparable. \textbf{Triage}
is a four-way classification over the disease classifier's structured output.
\textbf{Narrative generation} produces free text that no classifier can emit.

All figures reported below were obtained against a live, self-hosted
Qwen3-8B-AWQ teacher on RTX A6000 hardware. Where a figure is derived
arithmetically rather than measured, we state so explicitly.

\begin{table}[H]
\caption{Evaluation design. Sample size remains the principal limitation on the
narrative results, and we return to this in \S6.}
\begin{tabularx}{\textwidth}{@{}Yr L{56mm} l@{}}
\toprule
\textbf{Evaluation} & \textbf{n} & \textbf{Design} & \textbf{Strength} \\
\midrule
Triage, out-of-sample     & 160 & Unseen images, differing label distribution & \textbf{Strong} \\
Triage, cross-validation  & 80  & 5-fold $\times$ 5 seeds (25 fits)           & Strong \\
Triage, emergency stress  & 80  & Injected CRAO/AION cases                    & Synthetic, decisive \\
Narrative, variant sweep  & 24  & Six variants, single code path              & \textcolor{stopRed}{Limited} \\
Narrative, faithfulness   & 24  & Claim-level audit                           & \textcolor{stopRed}{Limited} \\
\bottomrule
\end{tabularx}
\end{table}

\section{Results: triage}

\subsection{Selected model}

We selected \code{feat\_logreg}, a multinomial logistic regression over 27
structured features derived from the classifier's output. The artefact is 3\,KB
of JSON weights and the forward pass is a single matrix multiplication, so the
serving path requires neither scikit-learn nor a deserialisation step at
start-up. We verified that the exported weights reproduce the fitted
estimator's predictions on all 240 available cases before adopting them.

\subsection{Comparison of architectures}

\begin{table}[H]
\caption{Candidate architectures on a common held-out split. Five of the seven
attained 1.000; we selected the smallest and fastest.}
\begin{tabularx}{\textwidth}{@{}Yrrrrrl@{}}
\toprule
\textbf{Model} & \textbf{Acc.} & \textbf{Macro-P} & \textbf{Macro-F1} &
\textbf{Size} & \textbf{p95} & \textbf{Gate} \\
\midrule
\rowcolor{softTeal}
\code{feat\_logreg}   & 1.000 & 1.000 & 1.000 & 3\,KB    & 0.14\,ms & Pass \\
\code{feat\_xgboost}  & 1.000 & 1.000 & 1.000 & 0.42\,MB & 1.8\,ms  & Pass \\
\code{rule\_baseline} & 1.000 & 1.000 & 1.000 & 0        & 0.02\,ms & Pass \\
\code{cnn\_triage}    & 0.792 & \textcolor{stopRed}{0.492} &
\textcolor{stopRed}{0.520} & 6.1\,MB & 161\,ms & \textcolor{stopRed}{Fail} \\
\bottomrule
\end{tabularx}
\end{table}

The image-based convolutional model underperformed for a diagnosable reason: it
infers triage from pixels, whereas the teacher derives it from the classifier's
structured output. Supplying that structured vector to a small tabular model
recovers the mapping directly.

\subsection{Generalisability}

We conducted four tests before producing any container image, in accordance with
the brief.

\begin{table}[H]
\begin{tabularx}{\textwidth}{@{}YL{54mm}@{}}
\toprule
\textbf{Test} & \textbf{Result} \\
\midrule
Out-of-sample (train 80, test 160 unseen images, differing label mix) &
Accuracy, macro-precision, F1 and $\kappa$ all \textbf{1.000} \\
Repeated stratified 5-fold cross-validation (25 fits) & \textbf{1.000 $\pm$ 0.000} \\
Learning curve & \textbf{1.000} from 20 training cases onward \\
EMERGENCY stress test (80 injected CRAO/AION) & Recall and precision \textbf{1.000} \\
\bottomrule
\end{tabularx}
\end{table}

\subsection{Safety provision}

No EMERGENCY case occurs anywhere in the training sample; consequently the
learned model cannot emit that class. Escalation for sight-threatening codes is
enforced by a deterministic override in the serving path, so a learned model
cannot downgrade such a finding. As no observed case exercises this path, we
cover it by unit test. The loader additionally refuses any artefact whose
feature layout has diverged from the one recorded at export, and degrades to the
deterministic rules on any failure rather than failing a screening.

\section{Results: narrative generation}

We distilled a 135M-parameter narrator and evaluated it at three precisions on
the held-out split. Output is plain prose rather than JSON, since the triage
fields are now supplied by the model in \S3; this removes the parse-failure mode
that rendered an earlier variant unusable, and generation succeeded in every
case at every precision.

\begin{table}[H]
\caption{Compact narrator, 24 held-out cases. Generation succeeds uniformly, so
fidelity rather than reliability distinguishes the variants.}
\begin{tabularx}{\textwidth}{@{}Yrrrrr@{}}
\toprule
\textbf{Precision} & \textbf{Size} & \textbf{Gen.\ rate} & \textbf{Omission} &
\textbf{Quoted probs} & \textbf{Acuity drift} \\
\midrule
\rowcolor{softTeal}
bf16 & 213.5\,MB & 1.000 & \textbf{0.000} & \textbf{0.000} & 0.250 \\
int8 & 107.3\,MB & \multicolumn{4}{l}{\textcolor{stopRed}{37.5\,s per case; evaluation impracticable}} \\
nf4  & 54.2\,MB  & 1.000 & \textcolor{stopRed}{\textbf{0.417}} &
\textcolor{stopRed}{0.167} & \textcolor{stopRed}{0.583} \\
\bottomrule
\end{tabularx}
\end{table}

\begin{callout}[stopRed]
\eyebrow{Grounds for withholding clinical enablement}

\vspace{3pt}
\textbf{The 4-bit variant substitutes a visible failure for an invisible one.}
It satisfies the 60\,MB constraint and produces fluent text in every case, while
omitting findings the teacher reported in 42\% of cases. Fluent omission is
considerably harder to detect in review than an absent output.

\vspace{4pt}
\textbf{The bf16 variant's clean scores do not extend beyond the evaluation
distribution.} The audit set is dominated by multi-finding cases resembling the
training data. Presented with a single-finding case constructed by hand, the
model reported ``a 75\% detection of diabetic retinopathy with a prevalence of
88\%'' for one finding at 88\% confidence; both figures are fabrications. A
single probe is not a measurement, but it demonstrates that the reported 0.000
is a property of the evaluation split rather than of the model.
\end{callout}

\subsection{Cost on the deployed hardware}

The deployed Space runs on free \code{cpu-basic} hardware (2 vCPU). Measured
there, restricted to two threads, the narrator requires 8.6\,s to load, 426.9\,MB
resident, and 6--10\,s per report, against a \code{/predict} endpoint that
currently responds in approximately 1\,s. We note additionally that
\code{transformers} is not a declared dependency of the deployed image, so
enabling the narrator would require an image rebuild, and that the smaller
quantised variants cannot run on this tier at all, since the bitsandbytes
kernels require CUDA. The narrator is therefore published but disabled by
default.

\section{Quantisation format study}

Our previous revision recommended AWQ or GPTQ over bitsandbytes, and projected
that GGUF \code{Q4\_K\_M} would reach approximately 80\,MB. We have since tested
the GGUF path and report that \textbf{the projection was incorrect}.

\begin{table}[H]
\caption{The AWQ/GPTQ row remains arithmetic, as neither library could be
installed in this environment. The GGUF rows are measured, using
\code{convert\_hf\_to\_gguf.py} and \code{llama-quantize} built from source.}
\begin{tabularx}{\textwidth}{@{}YL{28mm}L{34mm}rr@{}}
\toprule
\textbf{Format} & \textbf{Body} & \textbf{Embedding} & \textbf{BPW} & \textbf{Total} \\
\midrule
bitsandbytes NF4 & 4-bit & fp16, not quantised & --- & 109.8\,MB \\
bitsandbytes int8 & int8 & fp16, not quantised & --- & 107.3\,MB \\
AWQ / GPTQ (W4A16) & 4-bit & fp16, excluded by convention & --- & $\sim$110\,MB \\
\rowcolor{softTeal}
\textbf{GGUF \code{Q4\_K\_M}} & mixed & \textbf{q8\_0} & 6.17 & \textbf{105.5\,MB} \\
GGUF \code{Q5\_K\_M} & mixed & q8\_0 & --- & 112.1\,MB \\
GGUF \code{Q8\_0} & q8\_0 & q8\_0 & --- & 144.8\,MB \\
\bottomrule
\end{tabularx}
\end{table}

At 105.5\,MB the GGUF build is only marginally smaller than the 109.8\,MB
bitsandbytes build, rather than the substantial reduction our arithmetic
implied. We identified two architecture-specific causes and then tested the
explanation directly.

\begin{enumerate}[leftmargin=*,itemsep=3pt]
\item The embedding table \emph{is} quantised, confirming the principle we
argued, but to \code{q8\_0} rather than 4-bit: 54.00\,MiB to 28.69\,MiB, a factor
of 1.9 rather than the factor of 4 assumed.
\item The hidden size of SmolLM2-135M is 576, which is not divisible by 256 as
the \code{q4\_K} block structure requires. Consequently \textbf{180 tensors fell
back} to larger types, and llama.cpp reports the result as 6.17 bits per weight
against the 4.5 assumed.
\end{enumerate}

\subsection{Control experiment}

To establish that the divisibility constraint, rather than GGUF itself, accounts
for the discrepancy, we repeated the procedure on \code{pythia-160m}, whose
hidden size of 768 is divisible by 256. The comparison isolates the effect.

\begin{table}[H]
\caption{Control experiment. The only material difference between the two models
for this purpose is the hidden dimension.}
\begin{tabularx}{\textwidth}{@{}Yrr@{}}
\toprule
& \textbf{SmolLM2-135M (576)} & \textbf{pythia-160m (768)} \\
\midrule
\code{q4\_K} fallbacks        & 180        & \textbf{0} \\
Bits per weight               & 6.17       & \textbf{5.33} \\
Embedding quantised to        & \code{q8\_0} & \textbf{\code{q4\_K}} \\
Embedding reduction           & 1.9$\times$ & \textbf{3.6$\times$} \\
\bottomrule
\end{tabularx}
\end{table}

The control confirms the diagnosis: with a divisible hidden dimension the
fallback does not occur, the embedding receives 4-bit treatment, and the encoding
falls to 5.33 bits per weight. Applying that rate to our 134.5M-parameter
narrator gives approximately \textbf{89.6\,MB}, which would fall inside the
60--100\,MB band. We therefore conclude that the projection was directionally
sound but arithmetically optimistic, and that it fails for this model on account
of its width rather than on account of the format.

\subsection{Compatibility with vocabulary pruning}

We confirmed that the vocabulary-pruned checkpoint cannot be converted to GGUF:
the conversion script requires \code{vocab\_size} to equal the tokeniser size,
and the pruned model reports 920 against 49{,}152. Pruning and standard
quantisation formats are therefore alternative strategies rather than
complementary ones.

\subsection{Outstanding measurement}

Fidelity under GGUF remains unmeasured, as \code{llama-cli} did not return on
this host even for a trivial prompt. Whether \code{Q4\_K\_M} avoids the 42\%
omission observed in the pruned 4-bit build is therefore still open. We note
that at 105.5\,MB it would in any case compete against the bf16 build at
213.5\,MB rather than offer an edge-deployable option.

\section{Deployment and verification}

The platform is deployed to the Space \code{landwind22/Retinal\_Screening} and is
in service; \code{/health} reports \code{model\_loaded: true} across all 45
disease classes. It is the single deployment target, the previous Space having
been retired.

We record three defects in the deployment pipeline identified during this work,
all since corrected. The workflow previously reported success upon completion of
the git push, and on two occasions published a Space that subsequently failed to
start --- once through a README whose front matter was overwritten, and once
through a Dockerfile removed by \code{rsync --delete}. A third deployment was
lost when a push-triggered run and a manual dispatch executed concurrently and
raced to push. The workflow now polls the runtime stage and requires
\code{/health} to respond before reporting success, serialises runs, and retries
a rejected push. A six-hourly keep-alive maintains the Space against its 48-hour
inactivity timer.

\begin{table}[H]
\caption{Mobile audit (Lighthouse, Moto G Power emulation, throttled 4G),
representing the mid-range Android hardware this product targets.}
\begin{tabularx}{\textwidth}{@{}Yrrr@{}}
\toprule
& \textbf{Before} & \textbf{After} & \\
\midrule
Performance                 & 80      & \textbf{95}   & \\
Accessibility               & 96      & \textbf{100}  & \\
Largest Contentful Paint    & 5.0\,s  & \textbf{2.6\,s} & target 2.5\,s \\
Transferred payload         & 776\,KB & \textbf{293\,KB} & $-62$\% \\
Colour-contrast failures    & 31      & \textbf{0}    & \\
Redundant \code{alt} attributes & 2   & \textbf{0}    & \\
\bottomrule
\end{tabularx}
\end{table}

Two substantive defects were corrected. Text compression was not enabled, so
481\,KB of a 776\,KB payload was transmitted uncompressed at a cost of 2.4\,s on
a throttled connection; the remedy required \code{gzip\_proxied any}, since nginx
does not compress proxied responses by default and the payload in question is
proxied. Colour contrast failed on 31 nodes, of which the largest group arose not
from an unsuitable colour but from an opacity setting reducing a compliant
colour to a ratio of 2.71:1. We recomputed each value against its actual
background; we note that on dark sections the correction is lighter rather than
darker, a distinction that produced one regression before we identified it.

Container image sizes are 1.01\,GB compressed and 2.66\,GB on disk for the CPU
variant, and 5.90\,GB compressed for the GPU variant, of which 5.3\,GB is the
CUDA runtime and the CUDA build of PyTorch.

\section{Limitations}

We ask the supervisor to weigh the following when interpreting this report.

\begin{enumerate}[leftmargin=*,itemsep=3pt]
\item The narrative evaluations rest on 24 held-out cases and 56 training
examples. This is sufficient to separate a functioning model from a broken one
and little else; confidence intervals on the reported proportions are
correspondingly wide.
\item No clinician has reviewed the generated text. We regard this as the
principal barrier to clinical enablement.
\item No EMERGENCY case occurs in the 240 available traces, so escalation
behaviour is established only by injected stress testing.
\item The positive prediction path could not be exercised end-to-end against the
deployment, as the fundus images in the repository are dangling symbolic links
into a cleared cache. The rejection path was verified.
\item Lighthouse scores vary by approximately five points between runs on shared
infrastructure; single values should not be read as trends.
\end{enumerate}

\section{Recommendations and further work}

We recommend that the triage model be considered complete: it is validated,
computationally negligible, and removes a language-model dependency from the
decision path. We recommend that the grounded template remain the default
narrative and that \code{NARRATOR\_ENABLED} remain false, on the grounds that
the available instrumentation cannot yet distinguish a sound narrator from an
unsound one, which in a clinical screening context is the property that must be
established first.

We propose the following sequence of further work:

\begin{enumerate}[leftmargin=*,itemsep=3pt]
\item Regenerate the teacher traces at scale with the AI-disclosure statement
mandated, addressing the statistical-power and transparency limitations together.
The teacher is self-hosted, so the marginal cost is time alone.
\item Construct a faithfulness evaluation of appropriate rigour --- claim
decomposition with entailment scoring --- and obtain clinician review of a
sampled set.
\item Should an edge-deployable narrator remain desirable, evaluate a base model
whose hidden dimension is divisible by 256, as \S5.1 indicates this is the
condition under which the 60--100\,MB band becomes attainable without vocabulary
pruning.
\item Measure serving latency under vLLM before quoting any figure, as the
present measurements are dominated by framework overhead rather than by the
models.
\end{enumerate}

\vspace{14pt}
{\color{ruleGrey}\hrule height 0.6pt}
\vspace{6pt}

{\footnotesize\sffamily\color{inkMuted}\setlength{\parskip}{3pt}
\textbf{Supporting material.}\\
\code{docs/28-reasoner-cnn-vs-distilledqwen.md} --- design and results\\
\code{docs/29-narrator-verification-and-gaps.md} --- verification review\\
\code{outputs/reasoner\_comparison\_real/} \quad \code{outputs/generalizability/}

\textbf{Provenance.} Measurements were obtained on RTX A6000 hardware against a
live Qwen3-8B-AWQ teacher. Figures derived arithmetically rather than measured
are identified as such at each occurrence.\par}

\end{document}
